% ============================================================
% 3. MATERIALES Y MÉTODOS
% Una sección por fase del anteproyecto. Aquí va lo que se
% construyó y cómo se midió; los números van en el capítulo 4.
% Los comentarios "Repo:" señalan de dónde sale cada parte.
% ============================================================
\chapter{Materiales y métodos}
\label{cap:metodos}

\pendiente{un párrafo sobre la metodología en cascada y la validación en dos
niveles (referencias numéricas y hologramas de laboratorio), con todos los
métodos aplicados al mismo dato.}

% ------------------------------------------------------------
\section{Marco computacional unificado (Fase 1)}
\label{sec:implementacion}
% Repo: CamposT/propagadores.py, propagacion.py, backend.py (CPU/GPU),
%       campos.py; scripts/mpasm_minimo.py; código de referencia
%       en referencia/ y el parche de scripts/parchar_referencias.py.

\pendiente{arquitectura del paquete y la interfaz común de los tres
propagadores; soporte de GPU; cómo se reprodujo primero el código de
referencia \cite{zhao2020code}.}

\subsection{Verificación}
\label{sec:verificacion}
% Repo: CamposT/referencias.py, metricas.py; scripts/exactitud.py;
%       tests/test_rs1.py, test_propagadores.py.

\pendiente{referencia de Rayleigh--Sommerfeld en 1D; haz gaussiano analítico;
métrica SAM \cite{zhao2020}; pruebas automáticas.}

% ------------------------------------------------------------
\section{Calibración numérica (Fase 2)}
\label{sec:calibracion}
% Repo: scripts/comparacion.py, kf_evolucion.py, contraste_referencias.py;
%       resultados/kf, escalado_N.csv, escalado_s.csv.

\pendiente{parámetros barridos (distancias, paso y tamaño del sensor, $\lambda$,
$s$, $K_f$); cómo se construyen los mapas de error; criterio de «precisión
aceptable».}

\subsection{Iluminación esférica y distancia efectiva}
\label{sec:esferica}
% Repo: CamposT/montaje.py, retropropagacion.py; tarea 24 del cronograma.

\pendiente{el cambio de coordenadas antes de propagar: paso $\delta/M$ y la
distancia efectiva $z_\text{eq}$, validados contra RS.}

\subsection{Contraste con la transformada Z \emph{chirp}}
\label{sec:czt}

\pendiente{cómo se implementó la CZT y cómo se comparó con el MPASM.}

% ------------------------------------------------------------
\section{Montaje y registro experimental (Fase 3)}
\label{sec:montaje}
% Repo: docs/hoja_parametros_montaje.md; CamposT/montaje.py.

\pendiente{esquema del montaje; tabla de parámetros medidos ($\lambda$,
$\delta_x$, $\delta_y$, $N_x \times N_y$, $L$, $z$); muestras (USAF 1951 u
otra) y protocolo de adquisición.}

% ------------------------------------------------------------
\section{Reconstrucción y métricas de comparación (Fase 4)}
\label{sec:metricas}
% Repo: CamposT/pipeline.py, roi.py, metricas.py; scripts/barrido_foco.py,
%       retro_*.py, mi_prueba_*.py.

\pendiente{cadena de reconstrucción; búsqueda del plano de foco por nitidez
(sin verdad de terreno); resolución lateral sobre el \emph{target}, contraste,
relación señal--ruido y tiempo de cómputo.}
